\section{Background and Related Work}

In this section, we discuss why authors might wish to augment and enliven mathematical notation paired with dynamic representations, and then situate our system amid related work and tools. 

\subsection{Notation}
 
Mathematic notation is cognitively demanding: readers frequently shift attention between formulas and accompanying prose~\cite{Kohlhase2018-vv}, and readers perform worse on comprehension tasks when a text contains symbols rather than solely prose~\cite{Osterholm2006-al}. As a result, the notation students engage with often requires repeated exposure to understand~\cite{adams2003reading}. Notation has also been reported as an impediment to self-teaching technical topics like machine learning~\cite{Cai2019-lb}.

Subtle changes to the presentation of notation can improve its readability. Coloring and annotating formulas can reduce cognitive load in solving algebra problems~\cite{5c7148d6-681c-3534-885c-7936f89f5133}. Spacing between operators and operands affects how readers interpret order of operations~\cite{Landy2007-zk, Harrison2020-an}. Yet evidence suggest that a static augmentations is still received passively, and deeper understanding requires engagement that goes beyond reading. Chi and Wylie's ICAP framework~\cite{Chi2014-ks} provides empirical support for this intuition, showing that learning deepens as engagement moves from passive reception of information, to active manipulation of materials, to constructive generation of new inferences. A reader who can drag a variable and observe consequences is engaged actively and constructively — activating prior knowledge and inferring relationships. The possibility of deeper learning thus motivates the creation of explorable explanations.

\subsection{Explorable explanations and interactions}

The vision of interactive mathematical content on the web dates back to Bret Victor's ``Explorable Explanations''~\cite{victor2011explorable}, which argued that documents should let readers adjust assumptions and see consequences. Research on learning with multiple representations helps explain why this approach is effective. Ainsworth's DeFT framework~\cite{Ainsworth2006-pu} asserts that multiple representations serve three functions: they \textit{complement} each other, as when a formula supports precise specification while a graph supports perceptual pattern recognition; they \textit{constrain} interpretation, as when concrete visual behavior clarifies abstract symbolic relationships; and they help learners \textit{construct deeper understanding} when integrating across representations promotes abstraction and transfer. 

However, these benefits depend on learners being able to translate between representations---a task research consistently finds difficult. Tools that dynamically link formulas to visualizations could externalize this translation. Hutchins, Hollan, and Norman~\cite{Hutchins1985-wb} provide a cognitive account of how: when an interface lets users manipulate representations as if they were the objects themselves, users experience direct engagement. Interactive formulas aspire to have this quality---the notation becomes the manipulable object itself that could be linked to other representations. Head, Xie, and Hearst~\cite{Head2022-ws} provide the most comprehensive empirical study of formula augmentation practices through a content analysis of 47 documents containing 281 augmented formulas and 1,182 individual augmentations across 16 kinds. Notably, four kinds of interactive augmentations were observed: \textit{scrubbable expressions}, where readers click and drag on a value to change it in place; \textit{editable text fields}, where values can be typed into input boxes; \textit{linked selections}, where hovering over or clicking an expression highlights related content elsewhere in the document; and \textit{external controls}, where input widgets like sliders and drop-down menus outside the formula influence expression values.

To these four categories, the learning sciences suggest a fifth: \textit{formula walkthrough}. The segmenting principle states that learner comprehension improves when complex dynamic content is broken into segments whose advance the learner controls, rather than presented as a continuous sequence~\cite{Plass2009-sq, mayer2001multimedia}. De Jong and van Joolingen (1998) similarly found that gradually introducing complexity produced the strongest gains in learners' intuitive understanding. Dragunov and Herlocker~\cite{Dragunov2003-uy} proposed augmenting formulas with annotations showing how variables are manipulated across derivation stages, and controls that adjust the level of detail in a derivation.

Delta aims to make 5 augmentations - scrubbable expressions, editable text fields, linked selections, external controls, walkthroughs - all achievable with an intuitive grammar design.

\subsection{Tools for augmenting notation}

\paragraph{Markup languages.} LaTeX remains the dominant notation for static mathematical content. On the web, KaTeX~\cite{eisenberg2020katex} and MathJax~\cite{Cervone2012-fd} render LaTeX into high-quality static typeset output, but they are rendering engines that provide no affordances for authoring augmentations. Wu et al.~\cite{Wu2023-xc} introduced FFL, a CSS-like language for \textit{static} formula augmentation. Rather than interleaving styling commands within the LaTeX source, authors write selectors and properties in a companion style specification. FFL's study result showcased how a the separation of augmentation markup from formula markup can reduce viscosity and error proneness compared to a LaTeX baseline. However, FFL is limited to static augmentations. Taking a step towards interactivity, Crichton's Nota~\cite{crichton2021newmedium} and Li et al.'s HeartDown~\cite{Li2022-li} both let authors specify definitions of symbols, with HeartDown going further to make linear algebra papers executable, where editing formula values updates linked figures. While recognizing the importance of linked selections and multiple representations, neither addresses the breadth of basic interactions outlined by Head's empirical study \cite{Head2022-ws}.

\paragraph{Reading interfaces.} Research has explored augmentations that help readers engage with notation in context. Head et al.~\cite{Head2020-eg} designed ScholarPhi, an interface for scientific papers that allows readers to look up symbol definitions in tooltips and equation diagrams overlaid on display equations. Alcock and Wilkinson~\cite{warwick45724} designed e-Proofs, a system for presenting mathematical proofs where readers navigate a sequence of screens that selectively highlight parts of a proof and narrate relationships between lines. Augmented Math~\cite{Chulpongsatorn_2023} and Augmented Physics~\cite{Gunturu_2024} use machine learning and computer vision to retrofit interactive formulas onto static textbook pages. These systems demonstrate the demand for making explorable math notations.

\paragraph{Editing interfaces.} Existing editing interfaces for notation reveal a tension between structure and flexibility. WYSIWYG formula editors such as Microsoft Equation Editor~\cite{microsoft_equation_editor} and MathType~\cite{Topping1999UsingMT} preserve semantic structure but offer no affordances for interactivity. Authors who turn to vector graphics tools like Inkscape~\cite{inkscape} or Google Slides~\cite{google_slides} gain visual flexibility but produce output with no connection to the underlying mathematics --- change a value, and everything must be redrawn by hand. As one author in Head et al.'s study described, ``in Inkscape\ldots{} nothing's attached to anything\ldots{} everything's floating.'' Graspable 
Math~\cite{7821641} shows a third path: by keeping notation semantically grounded, it enables direct manipulation that reflects mathematical meaning, allowing learners to pinch expressions together or drag terms to rearrange equations. Yet its interactions are hardcoded to algebraic laws --- none of these tools let authors specify interactive behavior for their own notation.

\paragraph{Diagramming tools.} Powerful DSLs for mathematical diagramming have displayed a variety of language designs showcasing different strengths. Bluefish~\cite{10.1145/3654777.3676465} is a web-based diagramming framework that replaces React's component model with relations — building blocks like alignment, spacing, and arrows that can share children and specify layout. Grant Sanderson's manim~\cite{manim3b1b} library, as featured in his ~\cite{3blue1brown} video series, can produce rich mathematical animations in Python, and the output is a pre-rendered video. The language is shown to be more lower level in terms of controls, and Head et al.~\cite{Head2022-ws} report that one author estimated it would take a newcomer 1--2 months of use before understanding manim's conceptual model. FigureOne~\cite{figureone2018} is a JavaScript library for drawing, animating, and interacting with shapes, equations, and plots directly in the browser. Its flexible primitives support diagrams, animated equation derivations, and interactive figures. All of these tools feature primitives that are designed around the visual aspect of a mathematic diagram. Delta's design hopes to take a different approach by making use of the semantic and meaning encoded inside the mathematic structure. Adjacent to math, domain-specific primitives in a host language has proven effective in other technical domains. For instance, Hayatpur et al.'s Chisel~\cite{10.1145/3706598.3713723} is a DSL for data structure diagrams, grounded in a content analysis of 80 programmer-produced diagrams. Authors compose diagrams through declarative \textit{abstraction moves} --- selections, revisualizations, simplifications, and annotations --- that incrementally transform a concrete data display into a customized diagram. Another example, Penrose~\cite{10.1145/3386569.3392375}, translates mathematical notation into static diagrams via optimization. Bloom~\cite{Teller2025-pw} extends Penrose into the interactive domain: authors specify optimization constraints, and readers can drag diagram elements while the solver maintains those constraints in real time. Bloom also provides hooks for synchronizing internal diagram state with LaTeX on the page, enabling linked formula–diagram displays. Like Chisel, Delta aims to derives its primitives from an empirical analysis of existing artifacts specific to interactive mathematic explanations. Similar to Penrose, Delta hopes to recognize important semantic items that help govern the resulting artifact.

\subsection{Tools for interactive documents}

Students rely on surrounding prose text for interpretation ~\cite{Shepherd2014-gc} in most math articles. Thus, our work has a dimension of documents and embedding, so we draw on HCI research on interactive document affordances to guide Delta's design.

Existing explorations interactive documents shed light on the importance of \textit{synchronization}, where components are linked to shared states. Tangle.js~\cite{tangle2011} pioneered reactive documents with adjustable numbers: readers drag on numeric values in prose to adjust them, and dependent values update automatically. Inspired by inspired by Tangle.js, Curvenote developed ~\cite{curvenote_components, cockett2024continuous} a similar reactive web-component system. Conlen and Heer's Idyll~\cite{10.1145/3242587.3242600} introduced a reactive document markup language with a component-based architecture. Heer et al.'s Living Papers~\cite{Heer2023-ea} builds on Idyll's AST format to support articles with augmented equations, injecting annotations into KaTeX source and binding event listeners to rendered terms. 

In addition, there is a capability of \textit{substitution} in existing systems where values are substituted into variables in the state and used for computations. For instance, Litt et al.'s Potluck~\cite{potluck2022} demonstrates the broader vision of enriching text documents with reactive computation, by forming a loop of "extract data from text, compute with that data, and then display results back in the text." Similarly, Dragicevic et al.~\cite{10.1145/3290605.3300295} demonstrated ``explorable multiverse analyses,'' where readers can tinker with parameter choices in a research paper and see how results change, illustrating the power of parametric interaction within documents. 

Across all of these systems, a common design gap also emerges: reactivity is decoupled from mathematical structure such that none of the existing systems were able to addresses the basic interactions outlined by  Head's empirical study \cite{Head2022-ws}, to take a step back, none were able to have direct manipulation of formulas as a capability, the rendered formula remains a passive display surface that can be sometimes updated. Delta hopes to extend this space by providing the affordances created by existing systems while making formulas as first-class interactive objects.

\textbf{Computation and visualization platforms.} Several established platforms provide powerful computation and visualization for mathematics. Wolfram Mathematica~\cite{mathematica} offers deep computation and rich visualization, though its browser-based Wolfram Cloud has server-round-trip latency and a subscription requirement. Desmos~\cite{desmos} and GeoGebra~\cite{geogebra} are excellent embeddable tools for graphing and interactive geometry, respectively, but their embedded widgets remain self-contained calculator interfaces with their own sliders and input fields. Formulas within these platforms are not directly manipulable as interactive notation within a surrounding document. 

Across all of the tools surveyed in this section, a common limitation emerges: the rendered formula is not treated as a structured, interactive object. Reactivity operates at the variable-binding layer, entirely separate from the rendered notation. Computation platforms live in their own UI paradigms. None provide direct manipulation of formula symbols or formula-structure-aware interaction. An author who wants a reader to drag a variable in a rendered equation must manually wire overlay elements, event handlers, and state updates — a wide semantic distance, in Hutchins, Hollan, and Norman's~\cite{Hutchins1985-wb} terms, between the intention ("make $\alpha$ draggable") and the specification the system demands. This is exactly the kind of imperative entanglement that Delta seeks to eliminate.